% =====================================================================
%  Preámbulo compartido — Beamer
%  Curso: Análisis y Pronóstico de Datos Mineros con Python
%
%  Tema construido a mano sobre la base 'default' de Beamer:
%  compila con cualquier distribución TeX moderna sin paquetes de
%  tema externos (no requiere metropolis).
%
%  Cada deck de módulo hace:  % =====================================================================
%  Preámbulo compartido — Beamer
%  Curso: Análisis y Pronóstico de Datos Mineros con Python
%
%  Tema construido a mano sobre la base 'default' de Beamer:
%  compila con cualquier distribución TeX moderna sin paquetes de
%  tema externos (no requiere metropolis).
%
%  Cada deck de módulo hace:  % =====================================================================
%  Preámbulo compartido — Beamer
%  Curso: Análisis y Pronóstico de Datos Mineros con Python
%
%  Tema construido a mano sobre la base 'default' de Beamer:
%  compila con cualquier distribución TeX moderna sin paquetes de
%  tema externos (no requiere metropolis).
%
%  Cada deck de módulo hace:  \input{preambulo.tex}
% =====================================================================
\documentclass[11pt,aspectratio=169]{beamer}

\usepackage[T1]{fontenc}
\usepackage[utf8]{inputenc}
\usepackage[spanish,es-noshorthands]{babel}
\usepackage{lmodern}
\usepackage{amsmath,amssymb}
\usepackage{booktabs}
\usepackage{array}
\usepackage{xcolor}
\usepackage{graphicx}
\usepackage{listings}
\usepackage{tikz}
\usetikzlibrary{arrows.meta,positioning}

% ---------------------------------------------------------------------
%  Paleta
% ---------------------------------------------------------------------
\definecolor{azulmina}{HTML}{1F4E79}   % azul profundo
\definecolor{cobre}{HTML}{B5651D}      % acento cobre
\definecolor{griscarbon}{HTML}{2B2B2B}
\definecolor{grissuave}{HTML}{EDEEF0}
\definecolor{verdeok}{HTML}{2E7D32}
\definecolor{rojoalerta}{HTML}{C62828}

% ---------------------------------------------------------------------
%  Tema visual (sobre 'default')
% ---------------------------------------------------------------------
\usetheme{default}
\usefonttheme{professionalfonts}
\setbeamertemplate{navigation symbols}{}
\setbeamercolor{normal text}{fg=griscarbon}
\setbeamercolor{structure}{fg=azulmina}
\setbeamercolor{alerted text}{fg=rojoalerta}
\setbeamercolor{frametitle}{fg=azulmina}
\setbeamerfont{frametitle}{series=\bfseries,size=\large}
\setbeamertemplate{frametitle}{%
  \vskip6pt%
  \usebeamerfont{frametitle}\usebeamercolor[fg]{frametitle}\insertframetitle%
  \vskip2pt{\color{cobre}\hrule height 1.1pt}%
}
\setbeamertemplate{itemize items}[circle]
\setbeamertemplate{itemize subitem}[triangle]
\setbeamercolor{itemize item}{fg=cobre}
\setbeamercolor{itemize subitem}{fg=cobre}
\setbeamercolor{enumerate item}{fg=cobre}

\setbeamercolor{block title}{bg=azulmina,fg=white}
\setbeamercolor{block body}{bg=grissuave,fg=griscarbon}
\setbeamercolor{block title alerted}{bg=rojoalerta,fg=white}
\setbeamercolor{block title example}{bg=verdeok,fg=white}

% Barra inferior compacta
\setbeamertemplate{footline}{%
  \hbox{%
    \begin{beamercolorbox}[wd=.62\paperwidth,ht=2.4ex,dp=1ex,leftskip=1em]{author in head/foot}%
      \color{white}\scriptsize Análisis y Pronóstico de Datos Mineros con Python%
    \end{beamercolorbox}%
    \begin{beamercolorbox}[wd=.38\paperwidth,ht=2.4ex,dp=1ex,rightskip=1em]{title in head/foot}%
      \color{white}\scriptsize\hfill\insertshortsubtitle\ \ \textbullet\ \ \insertframenumber/\inserttotalframenumber%
    \end{beamercolorbox}}%
  \vskip0pt%
}
\setbeamercolor{author in head/foot}{bg=azulmina}
\setbeamercolor{title in head/foot}{bg=cobre}

% ---------------------------------------------------------------------
%  Código Python
% ---------------------------------------------------------------------
\lstdefinestyle{py}{
  language=Python,
  basicstyle=\ttfamily\scriptsize,
  keywordstyle=\color{azulmina}\bfseries,
  stringstyle=\color{cobre},
  commentstyle=\color{verdeok}\itshape,
  showstringspaces=false,
  breaklines=true,
  breakatwhitespace=true,
  columns=flexible,
  keepspaces=true,
  frame=leftline,
  framerule=1.4pt,
  rulecolor=\color{cobre},
  xleftmargin=1.2em,
  literate={á}{{\'a}}1 {é}{{\'e}}1 {í}{{\'i}}1 {ó}{{\'o}}1 {ú}{{\'u}}1
           {ñ}{{\~n}}1 {Á}{{\'A}}1 {É}{{\'E}}1 {Í}{{\'I}}1 {Ó}{{\'O}}1
           {Ú}{{\'U}}1 {Ñ}{{\~N}}1
}
\lstset{style=py}
\newcommand{\code}[1]{\texttt{\small #1}}

% ---------------------------------------------------------------------
%  Ayudas
% ---------------------------------------------------------------------
\newcommand{\idea}[1]{%
  \begin{center}\setlength{\fboxsep}{8pt}%
  \colorbox{grissuave}{\parbox{0.9\linewidth}{\centering\color{azulmina}\bfseries #1}}%
  \end{center}}
\newcommand{\ojo}[1]{{\color{rojoalerta}\textbf{Ojo:} #1}}
\newenvironment{recordar}{\begin{block}{Qué deberías recordar}}{\end{block}}

% flecha de flujo (segura en modo texto y modo matemático)
\newcommand{\fl}{\ensuremath{\;{\color{cobre}\rightarrow}\;}}

% cadena de flujo (roadmap): se escala para caber en el ancho de la diapositiva
\newcommand{\ruta}[1]{\par\vskip2pt\begin{center}%
  \resizebox{0.97\linewidth}{!}{\ensuremath{\displaystyle #1}}%
  \end{center}\vskip2pt}

% ---------------------------------------------------------------------
%  Metadatos
% ---------------------------------------------------------------------
\title[Datos Mineros con Python]{Análisis y Pronóstico de Datos Mineros con Python}
\author{}
\date{16 h \textbullet\ 8 clases \textbullet\ Google Colab}

% Portada estándar del módulo:  \portada{N}{Título del módulo}
\newcommand{\portada}[2]{%
  \subtitle[Módulo #1]{Módulo #1 \textemdash\ #2}%
  \begin{frame}[plain]
    \vfill
    {\color{azulmina}\rule{\linewidth}{2pt}}\\[0.4em]
    {\large\bfseries\color{cobre}Módulo #1}\\[0.3em]
    {\LARGE\bfseries\color{azulmina} #2}\\[0.6em]
    {\color{azulmina}\rule{\linewidth}{2pt}}\\[1.2em]
    \textit{Análisis y Pronóstico de Datos Mineros con Python}\\[0.2em]
    \footnotesize Apuntes + notebook de Google Colab
    \vfill
  \end{frame}}

% =====================================================================
\documentclass[11pt,aspectratio=169]{beamer}

\usepackage[T1]{fontenc}
\usepackage[utf8]{inputenc}
\usepackage[spanish,es-noshorthands]{babel}
\usepackage{lmodern}
\usepackage{amsmath,amssymb}
\usepackage{booktabs}
\usepackage{array}
\usepackage{xcolor}
\usepackage{graphicx}
\usepackage{listings}
\usepackage{tikz}
\usetikzlibrary{arrows.meta,positioning}

% ---------------------------------------------------------------------
%  Paleta
% ---------------------------------------------------------------------
\definecolor{azulmina}{HTML}{1F4E79}   % azul profundo
\definecolor{cobre}{HTML}{B5651D}      % acento cobre
\definecolor{griscarbon}{HTML}{2B2B2B}
\definecolor{grissuave}{HTML}{EDEEF0}
\definecolor{verdeok}{HTML}{2E7D32}
\definecolor{rojoalerta}{HTML}{C62828}

% ---------------------------------------------------------------------
%  Tema visual (sobre 'default')
% ---------------------------------------------------------------------
\usetheme{default}
\usefonttheme{professionalfonts}
\setbeamertemplate{navigation symbols}{}
\setbeamercolor{normal text}{fg=griscarbon}
\setbeamercolor{structure}{fg=azulmina}
\setbeamercolor{alerted text}{fg=rojoalerta}
\setbeamercolor{frametitle}{fg=azulmina}
\setbeamerfont{frametitle}{series=\bfseries,size=\large}
\setbeamertemplate{frametitle}{%
  \vskip6pt%
  \usebeamerfont{frametitle}\usebeamercolor[fg]{frametitle}\insertframetitle%
  \vskip2pt{\color{cobre}\hrule height 1.1pt}%
}
\setbeamertemplate{itemize items}[circle]
\setbeamertemplate{itemize subitem}[triangle]
\setbeamercolor{itemize item}{fg=cobre}
\setbeamercolor{itemize subitem}{fg=cobre}
\setbeamercolor{enumerate item}{fg=cobre}

\setbeamercolor{block title}{bg=azulmina,fg=white}
\setbeamercolor{block body}{bg=grissuave,fg=griscarbon}
\setbeamercolor{block title alerted}{bg=rojoalerta,fg=white}
\setbeamercolor{block title example}{bg=verdeok,fg=white}

% Barra inferior compacta
\setbeamertemplate{footline}{%
  \hbox{%
    \begin{beamercolorbox}[wd=.62\paperwidth,ht=2.4ex,dp=1ex,leftskip=1em]{author in head/foot}%
      \color{white}\scriptsize Análisis y Pronóstico de Datos Mineros con Python%
    \end{beamercolorbox}%
    \begin{beamercolorbox}[wd=.38\paperwidth,ht=2.4ex,dp=1ex,rightskip=1em]{title in head/foot}%
      \color{white}\scriptsize\hfill\insertshortsubtitle\ \ \textbullet\ \ \insertframenumber/\inserttotalframenumber%
    \end{beamercolorbox}}%
  \vskip0pt%
}
\setbeamercolor{author in head/foot}{bg=azulmina}
\setbeamercolor{title in head/foot}{bg=cobre}

% ---------------------------------------------------------------------
%  Código Python
% ---------------------------------------------------------------------
\lstdefinestyle{py}{
  language=Python,
  basicstyle=\ttfamily\scriptsize,
  keywordstyle=\color{azulmina}\bfseries,
  stringstyle=\color{cobre},
  commentstyle=\color{verdeok}\itshape,
  showstringspaces=false,
  breaklines=true,
  breakatwhitespace=true,
  columns=flexible,
  keepspaces=true,
  frame=leftline,
  framerule=1.4pt,
  rulecolor=\color{cobre},
  xleftmargin=1.2em,
  literate={á}{{\'a}}1 {é}{{\'e}}1 {í}{{\'i}}1 {ó}{{\'o}}1 {ú}{{\'u}}1
           {ñ}{{\~n}}1 {Á}{{\'A}}1 {É}{{\'E}}1 {Í}{{\'I}}1 {Ó}{{\'O}}1
           {Ú}{{\'U}}1 {Ñ}{{\~N}}1
}
\lstset{style=py}
\newcommand{\code}[1]{\texttt{\small #1}}

% ---------------------------------------------------------------------
%  Ayudas
% ---------------------------------------------------------------------
\newcommand{\idea}[1]{%
  \begin{center}\setlength{\fboxsep}{8pt}%
  \colorbox{grissuave}{\parbox{0.9\linewidth}{\centering\color{azulmina}\bfseries #1}}%
  \end{center}}
\newcommand{\ojo}[1]{{\color{rojoalerta}\textbf{Ojo:} #1}}
\newenvironment{recordar}{\begin{block}{Qué deberías recordar}}{\end{block}}

% flecha de flujo (segura en modo texto y modo matemático)
\newcommand{\fl}{\ensuremath{\;{\color{cobre}\rightarrow}\;}}

% cadena de flujo (roadmap): se escala para caber en el ancho de la diapositiva
\newcommand{\ruta}[1]{\par\vskip2pt\begin{center}%
  \resizebox{0.97\linewidth}{!}{\ensuremath{\displaystyle #1}}%
  \end{center}\vskip2pt}

% ---------------------------------------------------------------------
%  Metadatos
% ---------------------------------------------------------------------
\title[Datos Mineros con Python]{Análisis y Pronóstico de Datos Mineros con Python}
\author{}
\date{16 h \textbullet\ 8 clases \textbullet\ Google Colab}

% Portada estándar del módulo:  \portada{N}{Título del módulo}
\newcommand{\portada}[2]{%
  \subtitle[Módulo #1]{Módulo #1 \textemdash\ #2}%
  \begin{frame}[plain]
    \vfill
    {\color{azulmina}\rule{\linewidth}{2pt}}\\[0.4em]
    {\large\bfseries\color{cobre}Módulo #1}\\[0.3em]
    {\LARGE\bfseries\color{azulmina} #2}\\[0.6em]
    {\color{azulmina}\rule{\linewidth}{2pt}}\\[1.2em]
    \textit{Análisis y Pronóstico de Datos Mineros con Python}\\[0.2em]
    \footnotesize Apuntes + notebook de Google Colab
    \vfill
  \end{frame}}

% =====================================================================
\documentclass[11pt,aspectratio=169]{beamer}

\usepackage[T1]{fontenc}
\usepackage[utf8]{inputenc}
\usepackage[spanish,es-noshorthands]{babel}
\usepackage{lmodern}
\usepackage{amsmath,amssymb}
\usepackage{booktabs}
\usepackage{array}
\usepackage{xcolor}
\usepackage{graphicx}
\usepackage{listings}
\usepackage{tikz}
\usetikzlibrary{arrows.meta,positioning}

% ---------------------------------------------------------------------
%  Paleta
% ---------------------------------------------------------------------
\definecolor{azulmina}{HTML}{1F4E79}   % azul profundo
\definecolor{cobre}{HTML}{B5651D}      % acento cobre
\definecolor{griscarbon}{HTML}{2B2B2B}
\definecolor{grissuave}{HTML}{EDEEF0}
\definecolor{verdeok}{HTML}{2E7D32}
\definecolor{rojoalerta}{HTML}{C62828}

% ---------------------------------------------------------------------
%  Tema visual (sobre 'default')
% ---------------------------------------------------------------------
\usetheme{default}
\usefonttheme{professionalfonts}
\setbeamertemplate{navigation symbols}{}
\setbeamercolor{normal text}{fg=griscarbon}
\setbeamercolor{structure}{fg=azulmina}
\setbeamercolor{alerted text}{fg=rojoalerta}
\setbeamercolor{frametitle}{fg=azulmina}
\setbeamerfont{frametitle}{series=\bfseries,size=\large}
\setbeamertemplate{frametitle}{%
  \vskip6pt%
  \usebeamerfont{frametitle}\usebeamercolor[fg]{frametitle}\insertframetitle%
  \vskip2pt{\color{cobre}\hrule height 1.1pt}%
}
\setbeamertemplate{itemize items}[circle]
\setbeamertemplate{itemize subitem}[triangle]
\setbeamercolor{itemize item}{fg=cobre}
\setbeamercolor{itemize subitem}{fg=cobre}
\setbeamercolor{enumerate item}{fg=cobre}

\setbeamercolor{block title}{bg=azulmina,fg=white}
\setbeamercolor{block body}{bg=grissuave,fg=griscarbon}
\setbeamercolor{block title alerted}{bg=rojoalerta,fg=white}
\setbeamercolor{block title example}{bg=verdeok,fg=white}

% Barra inferior compacta
\setbeamertemplate{footline}{%
  \hbox{%
    \begin{beamercolorbox}[wd=.62\paperwidth,ht=2.4ex,dp=1ex,leftskip=1em]{author in head/foot}%
      \color{white}\scriptsize Análisis y Pronóstico de Datos Mineros con Python%
    \end{beamercolorbox}%
    \begin{beamercolorbox}[wd=.38\paperwidth,ht=2.4ex,dp=1ex,rightskip=1em]{title in head/foot}%
      \color{white}\scriptsize\hfill\insertshortsubtitle\ \ \textbullet\ \ \insertframenumber/\inserttotalframenumber%
    \end{beamercolorbox}}%
  \vskip0pt%
}
\setbeamercolor{author in head/foot}{bg=azulmina}
\setbeamercolor{title in head/foot}{bg=cobre}

% ---------------------------------------------------------------------
%  Código Python
% ---------------------------------------------------------------------
\lstdefinestyle{py}{
  language=Python,
  basicstyle=\ttfamily\scriptsize,
  keywordstyle=\color{azulmina}\bfseries,
  stringstyle=\color{cobre},
  commentstyle=\color{verdeok}\itshape,
  showstringspaces=false,
  breaklines=true,
  breakatwhitespace=true,
  columns=flexible,
  keepspaces=true,
  frame=leftline,
  framerule=1.4pt,
  rulecolor=\color{cobre},
  xleftmargin=1.2em,
  literate={á}{{\'a}}1 {é}{{\'e}}1 {í}{{\'i}}1 {ó}{{\'o}}1 {ú}{{\'u}}1
           {ñ}{{\~n}}1 {Á}{{\'A}}1 {É}{{\'E}}1 {Í}{{\'I}}1 {Ó}{{\'O}}1
           {Ú}{{\'U}}1 {Ñ}{{\~N}}1
}
\lstset{style=py}
\newcommand{\code}[1]{\texttt{\small #1}}

% ---------------------------------------------------------------------
%  Ayudas
% ---------------------------------------------------------------------
\newcommand{\idea}[1]{%
  \begin{center}\setlength{\fboxsep}{8pt}%
  \colorbox{grissuave}{\parbox{0.9\linewidth}{\centering\color{azulmina}\bfseries #1}}%
  \end{center}}
\newcommand{\ojo}[1]{{\color{rojoalerta}\textbf{Ojo:} #1}}
\newenvironment{recordar}{\begin{block}{Qué deberías recordar}}{\end{block}}

% flecha de flujo (segura en modo texto y modo matemático)
\newcommand{\fl}{\ensuremath{\;{\color{cobre}\rightarrow}\;}}

% cadena de flujo (roadmap): se escala para caber en el ancho de la diapositiva
\newcommand{\ruta}[1]{\par\vskip2pt\begin{center}%
  \resizebox{0.97\linewidth}{!}{\ensuremath{\displaystyle #1}}%
  \end{center}\vskip2pt}

% ---------------------------------------------------------------------
%  Metadatos
% ---------------------------------------------------------------------
\title[Datos Mineros con Python]{Análisis y Pronóstico de Datos Mineros con Python}
\author{}
\date{16 h \textbullet\ 8 clases \textbullet\ Google Colab}

% Portada estándar del módulo:  \portada{N}{Título del módulo}
\newcommand{\portada}[2]{%
  \subtitle[Módulo #1]{Módulo #1 \textemdash\ #2}%
  \begin{frame}[plain]
    \vfill
    {\color{azulmina}\rule{\linewidth}{2pt}}\\[0.4em]
    {\large\bfseries\color{cobre}Módulo #1}\\[0.3em]
    {\LARGE\bfseries\color{azulmina} #2}\\[0.6em]
    {\color{azulmina}\rule{\linewidth}{2pt}}\\[1.2em]
    \textit{Análisis y Pronóstico de Datos Mineros con Python}\\[0.2em]
    \footnotesize Apuntes + notebook de Google Colab
    \vfill
  \end{frame}}


\begin{document}
\portada{1}{Preparación de datos con Python}

% ---------------------------------------------------------------
\begin{frame}{Objetivo del módulo}
Tomar una base de un proceso industrial y dejarla lista para el análisis temporal:
\begin{itemize}
  \item cargarla correctamente (formato, separador, decimal, codificación);
  \item inspeccionar su estructura y sus tipos;
  \item corregir problemas básicos de calidad;
  \item organizar la variable temporal;
  \item hacer una primera exploración gráfica.
\end{itemize}
\vfill
\idea{Datos originales \fl datos limpios y estructurados}
\end{frame}

\begin{frame}{No se trata de aprender Python de memoria}
En el curso se trabaja con \textbf{notebooks preparados en Google Colab}. Lo importante:
\begin{itemize}
  \item reconocer la estructura del código;
  \item ejecutar las celdas \textbf{de arriba hacia abajo};
  \item identificar qué variables hay que modificar;
  \item interpretar las salidas;
  \item adaptar el análisis a una nueva base.
\end{itemize}
\vfill
\[
\text{carga}\fl\text{limpieza}\fl\text{análisis}\fl\text{modelo}
\]
Si una celda anterior no se ejecutó, las siguientes fallan.
\end{frame}

\begin{frame}[fragile]{Tres librerías durante las primeras sesiones}
\begin{columns}[T]
\column{0.33\linewidth}
\textbf{pandas}\\ datos tabulares: importar, filtrar, fechas, faltantes, agrupar.
\column{0.33\linewidth}
\textbf{NumPy}\\ operaciones numéricas y valores especiales (\code{np.nan}).
\column{0.33\linewidth}
\textbf{matplotlib}\\ gráficos: \code{plot}, \code{hist}, \code{scatter}, \code{boxplot}.
\end{columns}
\vfill
\begin{lstlisting}
import pandas as pd
import numpy as np
import matplotlib.pyplot as plt
\end{lstlisting}
No necesitamos cada librería completa: solo las herramientas del problema.
\end{frame}

\begin{frame}{La estructura de una base: el DataFrame}
\begin{center}
\begin{tabular}{lrrrr}
\toprule
Fecha & Tonelaje\_tph & Ley\_Cu & Recuperacion\_pct & Potencia\_kW\\
\midrule
01-08-2026 08:00 & 1450 & 0.82 & 88.2 & 3150\\
01-08-2026 09:00 & 1472 & 0.80 & 87.9 & 3180\\
01-08-2026 10:00 & 1435 & 0.85 & 89.1 & 3105\\
\bottomrule
\end{tabular}
\end{center}
\begin{itemize}
  \item cada \textbf{columna} es una variable; cada \textbf{fila}, una observación;
  \item variables típicas: tonelaje, flujo, presión, temperatura, ley, recuperación, potencia\ldots
  \item y una variable especial: \idea{el tiempo}
\end{itemize}
\end{frame}

\begin{frame}[fragile]{Importación: los detalles importan}
\begin{lstlisting}
df = pd.read_csv("datos.csv")            # caso simple
df = pd.read_csv("datos.csv", sep=";")   # export en español
df = pd.read_csv("datos.csv", sep=";", decimal=",")   # 1450,5
df = pd.read_csv("datos.csv", encoding="latin1")      # acentos, enie
df = pd.read_excel("datos.xlsx", sheet_name="Produccion")
\end{lstlisting}
\vfill
\ojo{nunca asumir que una columna es numérica solo porque \emph{se ve} numérica.}
Si el decimal es coma y no lo indicamos, pandas la lee como texto.
\end{frame}

\begin{frame}[fragile]{Primera inspección: auditar antes de analizar}
\begin{lstlisting}
df.head()        df.tail()        df.shape       # (8760, 8)
df.columns       df.dtypes        df.info()
df.describe()
\end{lstlisting}
\begin{itemize}
  \item \code{info()}: filas, columnas, no nulos, tipos;
  \item nombres con espacios: \code{df.columns = df.columns.str.strip()};
  \item \code{describe()} revela lo sospechoso: \code{Tonelaje mínimo = -999}.
\end{itemize}
\ojo{un \code{-999} puede ser un código de "sin medición". Todavía no lo elimines: primero entiende qué significa.}
\end{frame}

\begin{frame}[fragile]{El tiempo como variable}
Aunque \emph{veamos} fechas, para Python \code{Fecha} suele ser \code{object} (texto).
\begin{lstlisting}
df["Fecha"] = pd.to_datetime(df["Fecha"], dayfirst=True)
df = df.sort_values("Fecha")     # el orden contiene información
df = df.set_index("Fecha")       # DatetimeIndex
df.loc["2026-08-01"]             # selección por día
\end{lstlisting}
\vfill
\[
\boxed{\text{en una serie temporal, el orden de las observaciones no puede ignorarse}}
\]
\end{frame}

\begin{frame}[fragile]{Frecuencia de muestreo y huecos}
\begin{lstlisting}
df = df.asfreq("h")            # impone paso horario; los huecos aparecen como NaN
df.resample("D").mean()        # horario -> diario (media)
df.resample("D").sum()         # producción total diaria
\end{lstlisting}
\begin{itemize}
  \item \code{asfreq()}: cambia/impone la frecuencia \emph{sin} agregar;
  \item \code{resample()}: agrupa y necesita una regla (\code{mean}, \code{sum}, \code{max}\ldots);
  \item la agregación correcta \textbf{depende de la variable}.
\end{itemize}
\end{frame}

\begin{frame}[fragile]{Valores faltantes: limpiar $\neq$ inventar}
\begin{lstlisting}
df.isna().sum()
100 * df.isna().mean()
df.dropna()                       # eliminar
df["Tonelaje"].interpolate()      # interpolar
# ...o mantener el NaN y documentarlo
\end{lstlisting}
La decisión depende de: cantidad, duración del hueco, variable, origen y finalidad.
\idea{No es lo mismo interpolar una hora perdida que reconstruir una detención de planta}
\end{frame}

\begin{frame}[fragile]{Duplicados y timestamps repetidos}
\begin{lstlisting}
df.duplicated().sum()
df = df.drop_duplicates()
df.index.duplicated().sum()       # dos lecturas para la misma hora
\end{lstlisting}
\vspace{0.4em}
Un timestamp repetido puede ser: dos mediciones legítimas, un duplicado, dos sensores, o algo a promediar.
\ojo{no se decide solo con Python: primero hay que saber cómo se generaron los datos.}
\end{frame}

\begin{frame}{Valores inconsistentes}
Una observación puede existir y no ser \code{NaN}, y aun así ser incorrecta:
\begin{itemize}
  \item \code{Tonelaje = -999}, \code{Recuperación = 240\%}, \code{Temperatura = -500\,\textdegree C};
  \item un valor \textbf{constante} durante muchas horas $\Rightarrow$ sensor congelado;
  \item cambios de escala, de unidades, saltos por recalibración.
\end{itemize}
\vfill
\[
\text{dato improbable}\ \neq\ \text{dato incorrecto}
\]
Si eliminamos todo lo inusual, borramos justo los eventos que luego queremos detectar.
\end{frame}

\begin{frame}[fragile]{Reemplazo de códigos conocidos}
Si la documentación del sistema dice que \code{-999} significa "sin medición":
\begin{lstlisting}
df["Tonelaje"] = df["Tonelaje"].replace(-999, np.nan)
\end{lstlisting}
\vspace{0.6em}
Ahora Python sabe \textbf{explícitamente} que la información está ausente, en lugar de
tratar \code{-999} como un número real que distorsiona media y desviación.
\end{frame}

\begin{frame}{Los cuatro gráficos y sus preguntas}
\begin{center}
\begin{tabular}{ll}
\toprule
Gráfico & Pregunta principal\\
\midrule
Serie temporal & ¿Cómo evoluciona la variable?\\
Histograma & ¿Cómo se distribuyen sus valores?\\
Boxplot & ¿Dispersión y valores atípicos?\\
Dispersión & ¿Cómo se relacionan dos variables?\\
\bottomrule
\end{tabular}
\end{center}
\vfill
El histograma \textbf{elimina el orden temporal}: dos series muy distintas pueden tener
el mismo histograma. Esa idea será central al pasar a series temporales.
\end{frame}

\begin{frame}[fragile]{Flujo mínimo del Módulo 1}
\begin{lstlisting}
df = pd.read_csv("datos_proceso.csv", sep=";", decimal=",")
df.info(); df.describe()
df["Tonelaje"] = df["Tonelaje"].replace(-999, np.nan)
df["Fecha"] = pd.to_datetime(df["Fecha"], dayfirst=True)
df = df.sort_values("Fecha").drop_duplicates().set_index("Fecha")
df = df.asfreq("h")
df.isna().sum()
df["Tonelaje"].plot(figsize=(12, 4))
\end{lstlisting}
\end{frame}

\begin{frame}{Errores frecuentes}
\begin{itemize}
  \item asumir tipo numérico sin verificar \code{dtypes};
  \item modelar sin haber mirado la serie;
  \item eliminar automáticamente todo valor extraño;
  \item interpolar un hueco largo como si fuera una medición aislada;
  \item olvidar \code{sort\_values} antes de trabajar la serie.
\end{itemize}
\end{frame}

\begin{frame}
\begin{recordar}
\begin{enumerate}
  \item Siempre inspeccionar antes de analizar.
  \item Verificar los tipos de las variables.
  \item Una fecha en texto todavía no es una variable temporal.
  \item En series temporales el orden importa.
  \item Revisar frecuencia de muestreo y huecos.
  \item Faltante, duplicado y anómalo son problemas \emph{distintos}.
  \item No eliminar automáticamente lo extraño.
  \item Antes de modelar, visualizar.
\end{enumerate}
\end{recordar}
\vfill
\[
\boxed{\text{La calidad del modelo nunca supera la calidad de los datos}}
\]
\end{frame}

\begin{frame}{Cierre y transición}
El Módulo 1 termina con una base cargada, inspeccionada, con el tiempo ordenado,
problemas reconocidos y una visualización preliminar.
\vfill
\begin{center}
M1: \emph{¿mis datos están bien preparados?}\\[0.4em]
\fl\quad M2: \emph{¿qué información estadística puedo extraer?}
\end{center}
\end{frame}

\end{document}
